\chapter{من الذرّات إلى المجرّات: مقاييس الكون}\label{ch:g9:atoms-to-galaxies}

قبل تسع سنوات بدأ هذا الكتاب بحواسّك الخمس وبأرجل خنفساء. وهو
يُختَم برحلة لا تستطيع الحواسّ وحدها أن تقطعها: إحدى وأربعون
قوة من قوى العشرة، من قلب \omterm{def:g9:atoms-to-galaxies:atom}{الذرّة} إلى حافة الكون المرصود ---
وأنت، بدقّة لا بأس بها، في المنتصف. فصل أخير، وسرّ أخير مستحقّ
من نكتة قديمة، وسلّم تقف عليه لكل الفيزياء الآتية.

\section{نزولًا: إلى داخل المادة}

\begin{definition}[الذرّات]\label{def:g9:atoms-to-galaxies:atom}
\omterm{def:g7:states-of-matter:molecule}{جزيئات} سنتك السابعة مبنيّة هي نفسها: فكل \omterm{def:g7:states-of-matter:molecule}{جزيء} تجمّع من
\emph{ذرّات}\index{ذرّة} --- أبجدية الطبيعة، ونحو مئة نوع،
تهجّي تراكيبها كل مادة. \omterm{def:g7:states-of-matter:molecule}{فجزيء} الماء: ذرّتا هيدروجين وذرّة
أكسجين. وللذرّة بدورها معمار: \emph{نواة}\index{نواة} ضئيلة
ثقيلة موجبة الشحنة في المركز، وبعيدًا خارجها سحابة من
\emph{الإلكترونات}\index{إلكترون} --- حاملات شحنة سالبة تكاد
لا تزن شيئًا. وأحجام الذرّات قرب $10^{-10}$ \unit{m}؛ ونواها
قرب $10^{-15}$ --- أصغر بمئة ألف مرة ثانية.
\end{definition}

\begin{proposition}[المادة خواء في معظمها]\label{prop:g9:atoms-to-galaxies:empty}
كبّر \omterm{def:g9:atoms-to-galaxies:atom}{ذرّة} إلى ملعب عظيم، تصر نواتها حبّة بازلّاء عند دائرة
المنتصف --- \omterm{def:g9:atoms-to-galaxies:atom}{والإلكترونات} بعوضًا في أعلى المدرّجات: وكل ما
بينهما خالٍ. فطاولة الرخام، وسندان الحديد، ويدك أنت:
خواء في الأغلب الأعمّ، وصلابتها رفض كهربائي --- سحب إلكترونات
\omterm{def:g9:atoms-to-galaxies:atom}{الذرّات} تتنافر فتصنع وهم الامتلاء. وهي من أذهل جمل الفيزياء،
وهي حساب صريح للمقدار $10^{-10}$ في مقابل $10^{-15}$.
\end{proposition}

\begin{example}[النكتة القديمة، مدفوعةً]\label{ex:g9:atoms-to-galaxies:electron}
كانت فصول الكهرباء مدينة لك بنكتة صغيرة، وها هي. فالشحنات
السائرة في كل سلك في هذا الكتاب \emph{إلكترونات} --- وهي
سالبة، فتسير من مأخذ \omterm{def:g3:first-electric-circuit:battery}{البطارية} $-$ إلى مأخذها $+$: أي عكس
الاتجاه الاصطلاحي الذي ثبّته الفيزيائيون قبل قرن من أن يستطيع
أحد أن يسأل السائرين. وقد أُبقي الاصطلاح --- فكل قانون من
قوانين دواراتك يعمل معه عملًا تامًّا --- لكن المفارقة الصغيرة
قائمة: ففي كل مخطّط رسمته، يمشي الحشد الحقيقي في الجهة
الأخرى.
\end{example}

\begin{example}[وأعمق نزولًا]\label{ex:g9:atoms-to-galaxies:deeper}
وللنواة أجزاؤها هي، وحكايتها --- كيف ترتبط، ولماذا تنشطر بعض
النوى (\omterm{def:g5:heat-and-insulation:heat}{حرارة} محطّة التوليد) أو تندمج (\omterm{def:g5:heat-and-insulation:heat}{حرارة} الشمس) --- تنزل
بالفيزياء إلى $10^{-15}$ \omterm{def:g2:measuring-length:units}{متر} وما دونه، إلى أصغر البنى التي
سبرها البشر. ومجلّد الثانوية يفتح \omterm{def:g9:atoms-to-galaxies:atom}{النواة} كما ينبغي؛ وسنوات
الجامعة تمضي أبعد. فدرجات السلّم السفلى ما زالت تُبنى.
\end{example}

\section{صعودًا: إلى السماء}

\begin{example}[الصعود]\label{ex:g9:atoms-to-galaxies:ascent}
اصعد الآن، قوة عشرة بعد قوة، مارًّا بمعالم تسع سنوات المألوفة.
أنت، نحو $10^{0}$ \unit{m}. وساحة المدرسة، $10^{2}$. والأرض،
$10^{7}$ عرضًا --- كرة النموذج في الملعب الرياضي. وبعد \omterm{def:g2:moon-in-the-sky:moon}{القمر}،
$4 \times 10^{8}$؛ وبعد الشمس، $1.5 \times 10^{11}$ --- أي
ثماني دقائق ضوئية. وعرض النظام الشمسي، نحو $10^{13}$؛ وأقرب
\omterm{def:g4:solar-system:star}{نجم}، $4 \times 10^{16}$ --- أي أربع سنوات ضوئية. \omterm{ex:g5:stars-night-sky:milkyway}{ودرب
التبّانة}، $10^{21}$ مترًا من المدينة؛ والمجرّة المجاورة، على
بعد $2 \times 10^{22}$؛ وأعمق سماء مسحت، قرب $10^{26}$ مترًا
--- \omterm{def:g1:light-and-shadows:source}{وضوء} مجرّاتها أقدم من الأرض تحت كرسيّك.
\end{example}

\begin{omfigure}
\begin{tikzpicture}[scale=0.85]
  \draw[-Stealth, very thick] (0,0.6) -- (11.4,0.6);
  \foreach \x/\p in {0.4/-15, 1.7/-10, 3.0/-5, 4.3/0, 5.6/5,
      6.9/10, 8.2/15, 9.5/20, 10.8/25} {
    \draw[thick] (\x,0.45) -- (\x,0.75);
    \node[below, font=\footnotesize] at (\x,0.4) {$10^{\p}$};
  }
  % labels on two alternating rows, with thin leaders
  \foreach \x/\h/\name in {0.4/0.95/نواة, 1.7/1.6/ذرّة,
      3.0/0.95/نملة, 4.3/1.6/أنت, 5.6/0.95/{بلد},
      6.9/1.6/{بُعد الشمس}, 8.2/0.95/{أقرب النجوم},
      9.5/1.6/{المجرّة}, 10.8/0.95/{السماء العميقة}} {
    \draw[omExo!80, thin] (\x,0.75) -- (\x,\h);
    \node[above, font=\footnotesize] at (\x,\h) {\name};
  }
\end{tikzpicture}

{\small سلّم المقاييس، بالأمتار: إحدى وأربعون قوة من قوى
العشرة من \omterm{def:g9:atoms-to-galaxies:atom}{النواة} إلى السماء العميقة --- والقارئ واقف عند
الدرجة الوسطى تقريبًا. (وللنملة أن تسامحنا على التقريب.)}
\end{omfigure}

\begin{method}[التفكير برتب القدر]\label{met:g9:atoms-to-galaxies:oom}
أول أدوات الفيزيائي في أيّ سؤال جديد:
\begin{enumerate}
  \item عبّر عن المقدار بالكتابة العلمية؛
  \item واحتفظ بقوة العشرة وحدها --- \emph{رتبة
        القدر}\index{رتبة القدر}؛
  \item وقارن السلالم لا الأرقام: \omterm{def:g4:solar-system:star}{فنجم} عند $10^{16}$ \omterm{def:g2:measuring-length:units}{متر} في
        مقابل \omterm{def:g4:solar-system:planet}{كوكب} عند $10^{11}$ يختلفان بخمس درجات --- أي
        مئة ألف ضعف، مهما تقل الأرقام الأولى؛
  \item ثم اشحذ الأرقام بعد ذلك، إن لزم.
\end{enumerate}
تسع سنوات من تقدير الخطوات --- ``أهذا الجواب معقول؟'' ---
تنضج في هذه العادة. وأول فصل في مجلّد الثانوية يبني ثقافة
\omterm{def:g6:measurement-in-science:measuring}{القياس} عليها.
\end{method}

\begin{example}[حساب السلّم]\label{ex:g9:atoms-to-galaxies:arithmetic}
كم \omterm{def:g9:atoms-to-galaxies:atom}{ذرّة} تمتدّ على عرض قلم رصاص؟ القلم، $10^{-2}$ \unit{m}؛
\omterm{def:g9:atoms-to-galaxies:atom}{والذرّة}، $10^{-10}$: ثماني درجات --- أي مئة مليون \omterm{def:g9:atoms-to-galaxies:atom}{ذرّة}، كتفًا
إلى كتف. وكم أرضًا إلى الشمس؟ $1.5 \times 10^{11}$ على
$1.3 \times 10^{7}$: نحو $10^{4}$ --- أي عشرة آلاف كرة من
نموذج الملعب الرياضي، مصفوفة في صفّ. فالسلّم يجيب في ثوانٍ
عمّا تدفنه عدّات الكيلومترات الخام في الأصفار.
\end{example}

\section{المنظر من الدرجة الوسطى}

\begin{remark}[ما بنته تسع سنوات]\label{rem:g9:atoms-to-galaxies:built}
انظر إلى أسفل السلّم وأحصِ ما تملك. عند $10^{-10}$: \omterm{def:g7:states-of-matter:molecule}{الجزيئات}
\omterm{def:g9:atoms-to-galaxies:atom}{والذرّات} --- ومعها الحالات الثلاث، \omterm{def:g5:heat-and-insulation:heat}{والحرارة}، وتناقل \omterm{def:g3:sound-around-us:sound}{الصوت}،
وسائرو التيار. وحول $10^{0}$: القوى \omterm{def:g5:everyday-energy:energy}{والطاقة}، وأشعّة \omterm{def:g1:light-and-shadows:source}{الضوء}
المستقيمة، وقوانين الدارات --- فيزياء المقاس البشري في تسع
سنوات من التجارب. ومن $10^{7}$ إلى $10^{22}$: الأرض الدائرة،
والفصول المائلة، \omterm{def:g2:moon-in-the-sky:moon}{والقمر} الساقط، وقانون جذب واحد يدير البستان
والمجرّة سواءً. ولا مادة دراسية أخرى تسلّمك سلّمًا واحدًا من
داخل المادة إلى حافة المرئي --- وكل درجة صُعدت بآلات وأمانة
وحساب تستطيع أن تتحقّق منه بنفسك.
\end{remark}

\begin{remark}[ما ينتظر في الأعلى]\label{rem:g9:atoms-to-galaxies:ahead}
وأمّا العمل غير المنجَز فهو أطيب ما في الأمر. أيّ قانون يربط
\omterm{def:g2:pushes-and-pulls:force}{القوة} بالحركة بالضبط؟ وما \emph{هو} \omterm{def:g1:light-and-shadows:source}{الضوء} حتى يركض عند
$3 \times 10^{8}$ \omterm{def:g2:measuring-length:units}{متر} في الثانية ويحمل الألوان في داخله؟ وما
الذي ينشطر داخل \omterm{def:g9:atoms-to-galaxies:atom}{النواة} --- وما الذي يسبك \omterm{def:g1:light-and-shadows:source}{ضوء} الشمس؟ ولماذا
يجري إحصاء \omterm{def:g5:everyday-energy:energy}{الطاقة} النافعة إلى أسفل دائمًا؟ لقد أُثير كل سؤال
منها بأمانة وأُجّل بأمانة في هذه الصفحات؛ ولكل واحد جواب،
والأجوبة من أجمل ما تعرفه فصيلتنا. وهي تبدأ في مجلّد
الثانوية، بعد سنة من الآن. فاصحب السلّم --- وعادة التحقّق من
كل شيء.
\end{remark}

\section{تمارين}

\begin{exercise}[$\star$]\label{exo:g9:atoms-to-galaxies:1}
رتّب بحسب الحجم: \omterm{def:g7:states-of-matter:molecule}{جزيء}، \omterm{def:g9:atoms-to-galaxies:atom}{وذرّة}، \omterm{def:g9:atoms-to-galaxies:atom}{ونواة}، ونملة. وأعطِ لكل واحد
قوة العشرة الخاصّة به بالأمتار.
\end{exercise}

\begin{exercise}[$\star$]\label{exo:g9:atoms-to-galaxies:2}
صِف معمار \omterm{def:g9:atoms-to-galaxies:atom}{الذرّة} --- ما الذي يجلس في المركز، وما الذي يحيط
به، والحجمان اللذان يصنعان صورة الملعب.
\end{exercise}

\begin{exercise}[$\star$]\label{exo:g9:atoms-to-galaxies:3}
سدّد النكتة القديمة: من الذي يسير حقًّا في سلك، وفي أيّ جهة
--- ولماذا تنجو قوانين الدارات كلها من الكشف؟
\end{exercise}

\begin{exercise}[$\star$]\label{exo:g9:atoms-to-galaxies:4}
ما \omterm{met:g9:atoms-to-galaxies:oom}{رتبة القدر}؟ أعطِ رتبة قدر: طولك؛ وطول الصفّ؛ وقطر الأرض.
\end{exercise}

\begin{exercise}[$\star$]\label{exo:g9:atoms-to-galaxies:5}
بكم درجة من درجات السلّم تختلف هذه: \omterm{def:g9:atoms-to-galaxies:atom}{الذرّة} \omterm{def:g9:atoms-to-galaxies:atom}{والنواة}؟ وأنت
والأرض؟ وبُعد الشمس وبُعد أقرب \omterm{def:g4:solar-system:star}{نجم}؟
\end{exercise}

\begin{exercise}[$\star$]\label{exo:g9:atoms-to-galaxies:6}
``الطاولة خواء في معظمها.'' دافع عن الجملة بالعددين الذرّيين
--- واشرح ما الذي يجعل الطاولة تُحسّ صلبة مع ذلك.
\end{exercise}

\begin{exercise}[$\star$]\label{exo:g9:atoms-to-galaxies:7}
حساب السلّم: كم \omterm{def:g9:atoms-to-galaxies:atom}{ذرّة} تقريبًا، مصفوفةً في صفّ، تمتدّ على
ملّيمتر؟
\end{exercise}

\begin{exercise}[$\star$]\label{exo:g9:atoms-to-galaxies:8}
حدّد على السلّم، بقوة عشرة لكل واحد: \omterm{def:g8:speed-of-light:lightyear}{السنة الضوئية}؛ \omterm{ex:g5:stars-night-sky:milkyway}{ودرب
التبّانة}؛ وأعمق سماء مسحت.
\end{exercise}

\begin{exercise}[$\star\star$]\label{exo:g9:atoms-to-galaxies:9}
نموذج الملعب: إن كانت \omterm{def:g9:atoms-to-galaxies:atom}{النواة} حبّة بازلّاء قطرها \qty{1}{cm}،
فكم يبلغ عرض ملعب \omterm{def:g9:atoms-to-galaxies:atom}{الذرّة}؟ (النسبة $10^{5}$ --- حوّل إلى
الأمتار واحكم في مقابل ملعب حقيقي.)
\end{exercise}

\begin{exercise}[$\star\star$]\label{exo:g9:atoms-to-galaxies:10}
قطرة ماء تحوي نحو $10^{21}$ \omterm{def:g7:states-of-matter:molecule}{جزيء}. باستعمال مقارنة بحر القطرات
في فصل سنتك السابعة --- أو تقديرك أنت لعدد القطرات في
المحيطات كلها (نحو $10^{21}$ \omterm{def:g6:mass-and-volume:volume}{لتر}، وفي كل \omterm{def:g6:mass-and-volume:volume}{لتر} $10^{5}$ قطرة)
--- زِن الزعم القديم أن القطرة تحوي \omterm{def:g7:states-of-matter:molecule}{جزيئات} أكثر مما تحوي
البحار قطرات.
\end{exercise}

\begin{exercise}[$\star\star$]\label{exo:g9:atoms-to-galaxies:11}
جسمك يقف قرب منتصف السلّم: فكم درجة تقريبًا نزولًا إلى \omterm{def:g9:atoms-to-galaxies:atom}{الذرّة}،
وكم صعودًا إلى السماء العميقة؟ وعلامَ يدلّ هذا التناظر في شأن
مدى تسع سنوات من الدراسة؟
\end{exercise}

\begin{exercise}[$\star\star\star$]\label{exo:g9:atoms-to-galaxies:12}
تمرين الكتاب الأخير. اختر أيّ ثلاثة فصول من سنواتك التسع ---
واحدًا من الصفوف 1--3، وواحدًا من 4--6، وواحدًا من 7--9 ---
واكتب لكل واحد الجملة الوحيدة التي كنت ستقولها لنفسك الأصغر
وهي تبدأ ذلك الفصل: ماذا علّم حقًّا، منظورًا إليه من قمّة
السلّم. (ولا صفحة حلول لرسالة إلى النفس --- لكن اكتبها؛
فالمنهج القادم سيسرّه أنك فعلت.)
\end{exercise}

\section{مسألة: الرحلة بقوى العشرة}

\begin{problem}[{مسألة نهاية الأسبوع --- الصفّ يصوّر ``الرحلة
بقوى العشرة''؛ تكبير واحد، وإحدى وأربعون درجة؛ مؤتمر
النصّ}]\label{pb:g9:atoms-to-galaxies:1}
فيلم الصفّ: تكبير واحد متّصل من بساط نزهة إلى الخارج نحو
السماء العميقة، ثم إلى الداخل نحو \omterm{def:g9:atoms-to-galaxies:atom}{النواة} --- قوة عشرة واحدة
لكل ثانية من الفيلم. وأنت تكتب النصّ العلمي.

\medskip\noindent\textbf{القسم الأول --- إلى الخارج.}
يبدأ التكبير بتأطير \omterm{def:g2:measuring-length:units}{متر} واحد من البساط: $10^{0}$.
\begin{enumerate}
  \item عند أيّ ثوانٍ يحوي الإطار أول مرة: ساحة المدرسة
        كلها ($10^{2}$)؛ والأرض كلها ($10^{7}$)؛ \omterm{def:g6:sun-earth-moon:axisorbit}{ومدار}
        \omterm{def:g2:moon-in-the-sky:moon}{القمر} ($10^{9}$)؟
  \item وأيّ ثانية تؤطّر بُعد الشمس --- وماذا يقول المعلّق عن
        عمر \omterm{def:g1:light-and-shadows:source}{ضوء} الشمس، من فصل \omterm{def:g1:light-and-shadows:source}{الضوء}؟
  \item وبين أيّ ثانيتين يعبر الإطار الخواء العظيم حيث ينتهي
        النظام الشمسي ولم يدخل أقرب \omterm{def:g4:solar-system:star}{نجم} بعد؟ وماذا ينبغي أن
        تُظهر الشاشة هناك بأمانة؟
  \item وينتهي التكبير عند الثانية $26$: فما الذي يملأ
        الإطار، وكم يبلغ عمر أقدم \omterm{def:g1:light-and-shadows:source}{ضوء} فيه \omterm{def:g6:measurement-in-science:measuring}{بالقياس} إلى
        الأرض؟
\end{enumerate}

\medskip\noindent\textbf{القسم الثاني --- إلى الداخل.}
ينعكس التكبير فيغوص في ورقة على البساط.
\begin{enumerate}[resume]
  \item عند أيّ ثوانٍ يمرّ الفيلم بهذه: نملة ($10^{-2}$،
        بسخاء)؛ وخلية ($10^{-5}$)؛ \omterm{def:g7:states-of-matter:molecule}{وجزيء} ($10^{-9}$)؛
        \omterm{def:g9:atoms-to-galaxies:atom}{وذرّة} ($10^{-10}$)؟
  \item وبين سحابة إلكترونات \omterm{def:g9:atoms-to-galaxies:atom}{الذرّة} ونواتها، على الفيلم أن
        يكبّر خمس ثوانٍ أخرى عبر --- ماذا؟ وماذا يقول
        المعلّق عن الملعب؟
  \item وإطار الفيلم الأخير، عند $10^{-15}$: ما موضوعه ---
        وما سطر المعلّق الختامي الأمين عن الدرجات التي
        تحته؟
  \item واجمع طول الفيلم: كم ثانية إلى الخارج، وكم إلى
        الداخل، من \omterm{def:g2:measuring-length:units}{متر} البساط؟
\end{enumerate}

\medskip\noindent\textbf{القسم الثالث --- مؤتمر النصّ.}
\begin{enumerate}[resume]
  \item يعترض معلّم أن الفيلم ينفق ثانية على الخطوة من
        البيت إلى الشارع وثانية على الخطوة من مجرّة إلى
        عنقود مجرّات --- ``خطوات متفاوتة تفاوتًا صارخًا!''
        دافع عن تساوي الثواني بمنطق السلّم.
  \item ويريد المعلّق سطرًا متكرّرًا واحدًا للاتجاهين ---
        شيئًا صادقًا عند كل درجة عن الخواء. صُغه (فعلى
        النظام الشمسي \omterm{def:g9:atoms-to-galaxies:atom}{والذرّة} أن يندرجا تحته معًا).
  \item وعلى الشارة الختامية أن تسمّي الآلتين اللتين فتحتا
        الاتجاهين. سمِّهما، وسمِّ فصول هذا الكتاب التي
        قدّمت مبدأيهما.
  \item وكلمات الفيلم الأخيرة ملك المنهج نفسه: اكتب الخاتمة
        في جملتين --- تسع سنوات، وسلّم واحد، وما يبدأ في
        السنة القادمة.
\end{enumerate}
\end{problem}
